\chapter{Introduction}
\label{ch:intro}

\section{Motivation}

Large language models are increasingly deployed as autonomous agents: given an
operational goal---reduce claim-processing time, raise engagement, close support
tickets---they choose and take actions with limited human oversight. Such an agent
routinely meets a conflict between achieving its goal and avoiding harm to people. The
dangerous case is the \emph{harmful-but-effective} action, which advances the stated goal
while hurting a human stakeholder: a claims agent that hits its speed target by silently
denying borderline-valid claims, a moderation agent that raises engagement by amplifying
outrage. Two failure modes bracket the behaviour we want. An \emph{over-pragmatic} agent
takes the harmful shortcut to hit the goal; an \emph{over-cautious} agent refuses or
defers even when the effective action is perfectly safe, destroying its usefulness. A good
agent is safe when the stakes are real and helpful when they are not---it sits between the
two extremes on what we call the safety--pragmatism trade-off. ManagerBench
\citep{simhi2025managerbench} measures exactly where a model sits on this trade-off,
through two rates: Harm Avoidance (HA), the fraction of genuinely harmful dilemmas on
which the model takes the safe option, and Control Pragmatism (CP), the fraction of
harmless control items on which it still takes the effective action.

\section{Problem statement}

The model's weights are fixed---in the common case, a hosted model that cannot be
retrained. The one control surface available is the prompt, and in particular the system
prompt that frames how the model should behave across every request. This thesis asks:

\begin{quote}
Can we move a fixed LLM to a chosen point on its safety--pragmatism trade-off, cheaply and
interpretably, purely by optimizing the prompt---and can we build a tool that, given a
\emph{requested} safety and usefulness level, returns a prompt that achieves it, with a
statistical guarantee or an honest signal that the request is unreachable?
\end{quote}

Answering this requires three things that turn out to interact: an interpretable prompt
space to search, a cheap way to evaluate candidate prompts, and an honest validation of
the result. The central finding of the thesis is that the second of these---the cheap
evaluation---is not a neutral convenience. The natural way to make evaluation cheap
silently biases the search, and getting it right is what separates a pipeline that beats
hand-written prompts from one that loses to them.

\section{Contributions}

\begin{enumerate}
\item \textbf{A method.} Prompt design is cast as black-box optimization over two
interpretable continuous knobs---a safety-framing weight and a goal-pressure level---decoded
to a system prompt by a fixed template, scored on a cheap benchmark subset, and searched by
Bayesian optimization. The decode discretizes the space into 45 distinct prompts, which we
exploit for exact ground truth.
\item \textbf{The proxy-bias result (the core).} The intuitive ``keep the most informative
items'' subset---and the decision-boundary selections used by recent performance-guided
methods \citep{ipomp2025}---is a \emph{biased} estimator of the full benchmark; the bias is
configuration-conditional, cannot be removed by calibration, and, sitting inside the search
loop, steers Bayesian optimization to prompts that lose to the best hand-written prompt on
all four models. Representative sampling of the same budget is nearly unbiased.
\item \textbf{A repaired pipeline.} Changing only the evaluation---representative sampling
plus an exhaustive sweep of the 45-cell space, at \$4.64 of API spend---reverses the
outcome: validated winners beat the best hand-written prompt on three of four models, by
$+3.9$ to $+16.7$ MB points, all in a region no baseline had sampled.
\item \textbf{A characterization.} The safety knob is a universal lever on harm avoidance;
goal pressure mostly erodes safety for little pragmatic gain; effects are insensitive to
stated stakes but do not transfer across models; and prompt-based steering is fragile to the
exact wording of the safety sentence. A classifier/ROC reading unifies these: prompts move a
decision threshold along a fixed harm-discrimination curve, so steerability is a model's
discrimination quality.
\item \textbf{A controller.} Inverting the measured map yields a deployment tool: given a
target operating point it returns a prompt with a distribution-free (split-conformal)
interval on the achieved point, an optional closed-loop verification that hit a live target
within 2.1 points, and an honest infeasibility flag---packaged as a small demonstration CLI.
\end{enumerate}

\section{Roadmap}

Chapter~\ref{ch:background} gives the background. Chapter~\ref{ch:method} describes the
method and infrastructure. Chapters~\ref{ch:proxybias} and~\ref{ch:repair} are the core: the
proxy-bias result and the search damage it causes, with the repair. Chapter~\ref{ch:char}
characterizes how the knobs move each model. Chapter~\ref{ch:controller} presents the
controller. Chapter~\ref{ch:related} places the work in the literature,
Chapter~\ref{ch:discussion} discusses limitations and future work, and
Chapter~\ref{ch:conclusion} concludes.
